\section{Citações, notas e referências}\label{sec:citacoes-referencias}

Citações conectam o argumento do autor às fontes utilizadas; referências identificam os documentos citados ou utilizados de acordo com a política bibliográfica adotada. O template separa a apresentação das citações no texto da composição final da lista de referências.

\subsection{Citação indireta}

Na citação indireta, a ideia da fonte é apresentada com redação própria. Por exemplo, a separação entre conteúdo e apresentação é uma característica central do fluxo de preparação de documentos em \LaTeX{} \parencite{lamport1994}.

\guianormativa{ABNT NBR 10520:2023}{A indicação de autoria e data deve seguir o sistema de citação adotado no documento. O projeto usa \texttt{biblatex-abnt} como infraestrutura e mantém regressões específicas para a apresentação prevista no escopo testado \parencite{abnt10520}.}

\subsection{Citação direta curta}

Citações diretas curtas permanecem integradas ao parágrafo e usam sinais de citação apropriados. Para evitar reproduzir conteúdo protegido de uma norma técnica, o exemplo abaixo é deliberadamente sintético.

\guiaexemplo{Uma citação curta pode ser demonstrada como \enquote{texto citado com poucas linhas}, seguida da indicação bibliográfica correspondente. Em um trabalho real, o conteúdo entre aspas deve reproduzir fielmente a fonte consultada e trazer os dados exigidos para localização.}

\subsection{Citação direta longa}

\guianormativa{ABNT NBR 10520:2023 e escopo UFC auditado}{Citação direta com mais de três linhas é apresentada em bloco distinto, sem aspas, com recuo esquerdo de 4 cm, fonte reduzida de 10 pt e espaçamento simples \parencite{abnt10520}.}

O documento de referência exercita fisicamente esse comportamento com um texto sintético na Seção~\ref{sec:catalogo-exemplos}. A suíte mede recuo, tamanho de fonte, espaçamento, ausência de aspas e separação do bloco em relação ao corpo. Em uma citação direta real, a indicação bibliográfica deve incluir a página ou outro localizador aplicável quando a fonte o fornecer; o texto sintético do corpus geométrico não recebe um localizador fictício.

\subsection{Citação de citação}

O recurso conhecido como citação de citação deve ser usado com parcimônia, preferindo-se a consulta à fonte original sempre que possível. Quando necessário, a apresentação de \textit{apud} segue a política bibliográfica e a NBR 10520 vigente.

\guiapolitica{O repositório possui cenários específicos para verificar a apresentação de \textit{apud}. O fato de o template suportar o recurso não elimina a responsabilidade acadêmica de consultar a fonte original quando ela estiver disponível.}

\subsection{Notas de rodapé}

Notas podem complementar uma informação sem interromper excessivamente o fluxo do parágrafo. Sua tipografia reduzida foi descrita na Seção~\ref{sec:formatacao-geral}.

\guiaexemplo{Use notas para informação realmente complementar. Referências bibliográficas não devem ser transferidas arbitrariamente para notas apenas para contornar o sistema de citações configurado no trabalho.}

\subsection{Lista de referências}

\guianormativa{ABNT NBR 6023:2025}{A elaboração das referências segue a edição vigente adotada pelo projeto. No escopo de apresentação atualmente exercitado, as entradas são alinhadas à margem esquerda, usam fonte 12 pt e espaçamento simples internamente, com separação entre referências correspondente a uma linha em branco de espaço simples \parencite{abnt6023}.}

O arquivo \texttt{backmatter/references.bib} demonstra tipos bibliográficos diferentes: livro, artigo de periódico com DOI, trabalho em evento, dissertação, relatório técnico, norma e recurso eletrônico. O objetivo é exercitar os campos mais comuns sem sugerir que todos eles devam aparecer em todo trabalho.

\guiaexemplo{Uma referência com DOI só apresenta DOI quando esse identificador pertence ao documento referenciado. De modo semelhante, URL e data de acesso são usados quando pertinentes ao tipo de fonte e à regra bibliográfica aplicável.}

\subsection{Normas técnicas como referências}

As normas citadas neste guia aparecem também nas referências para que o leitor saiba qual edição orienta cada seção. O projeto não redistribui o texto integral das normas ABNT; a comunidade UFC pode utilizar os meios institucionais autorizados para consulta das normas técnicas.

\guiainstitucional{O Sistema de Bibliotecas da UFC disponibiliza orientação sobre acesso à coleção de normas técnicas. Este PDF identifica a norma e a edição adotadas, mas não substitui a consulta à publicação oficial \parencite{ufcnormastecnicas}.}

\index{citação direta}\index{citação indireta}\index{referências}\index{apud}\index{DOI}
